\section{Gap motion from the phase current}\label{sec:current}
The upper-anchor curvature theorem controls a time interval by a single complex tilt. A different identity computes the instantaneous gap velocity from the spatial curvature of a real current. It also determines where this phase coordinate can be used.

\subsection{The current and its domain}
On a real interval with \(B\ne0\), take a continuous argument \(\Theta\) of \(B\) and write \(B=e^{p+\ii\Theta}\). At the origin \(\Theta=0\). Define the current and its normalization by
\begin{equation}
 J(t,x)=\Imag(\overline{B(t,x)}B_x(t,x))=|B(t,x)|^2\Theta_x(t,x),
 \qquad j(t,x)=J(t,x)/(Z(t)^2a_t).
 \label{eq:current}
\end{equation}
The complete source determines the current even where the real part \(H\) vanishes. A positive current on \([0,X]\) gives a nonvanishing \(B\) and strictly increasing phase there. To include a labelled edge one must also reach its upper level \(q_{k+1}\).

For this particular lift one can identify a whole time range where the
coordinate fails between two positive real zeros. Put
\(t_{\rm m}=-\Phi''(0)/(2\Phi(0))>5/2\), where the last inequality
uses \(\Phi''(0)<-10\) and \(0<\Phi(0)<2\) \cite{EZPhaseScope}.
By \href{https://doi.org/10.1007/BF02075457}{Csordas and Varga, Theorem 2.1},
\(v\mapsto\log\Phi(\sqrt v)\) is strictly concave and its right derivative at
zero is \(-t_{\rm m}\). Hence \(f_t(u)=e^{tu^2}\Phi(u)\) is strictly
decreasing for \(u>0\) when \(t\le t_{\rm m}\). For \(x>0\), put \(h=\pi/x\).
Pairing half-periods gives
\[
 \Imag B(t,x)=\sum_{j\ge0}\int_0^h
 \{f_t(2jh+u)-f_t((2j+1)h+u)\}\sin(xu)\dd u>0 .
\]
The upper-half-plane argument consequently takes the value \(\pi/2\) at
both real zeros. It is nonconstant between them by analyticity of \(H\)
and \(H(t,0)>0\). Thus \(J\) takes both signs between the endpoints.
At any fixed time the same obstruction applies to sufficiently high positive
zeros, since integration by parts gives
\(\Imag B(t,x)=\Phi(0)/x+O_t(x^{-2})\).
The argument concerns the half-line lift. The exterior calculation in
Section~\ref{sec:exterior} uses its own zero counts and separation bounds.

\begin{theorem}[Velocity, interaction and joint forcing]\label{thm:velocity}
On a common increasing phase chart define \(v=(\log j)_x\) and
\(\mathcal A=(v(t,x_+)-v(t,x_-))/g\). Then
\begin{align}
 \dot x_\pm&=v(t,x_\pm),&
 (\log g)^{\displaystyle\cdot}&=\mathcal A
   =\frac1g\int_{x_-}^{x_+}(\log j)_{xx}\dd x, \label{eq:velocity}\\
 M&=\frac1{2D}-\frac{\mathcal A}{2},&
 \mathcal F&=\frac1{2D^2}+\frac12(\dot{\mathcal A}+\mathcal A^2).
 \label{eq:forcing}
\end{align}
The same \(v\) advances every edge in the chart. In the parent's normalization \(V=1-2DM=D\mathcal A\) and \(K=2D^2\mathcal F=1+D^2\ddot g/g\).
\end{theorem}
\begin{proof}
The imaginary part of \(B_t/B=-B_{xx}/B\) gives
\(\Theta_t=-\Theta_{xx}-2p_x\Theta_x\). Differentiating
\(\Theta(t,x_\pm(t))=q_{k},q_{k+1}\), respectively, yields
\(\dot x=2p_x+\Theta_{xx}/\Theta_x=(\log J)_x\). Its difference proves \eqref{eq:velocity}. Thus \(\dot D=2D\mathcal A\); substituting in both equations of \eqref{eq:parent} gives \eqref{eq:forcing}, including the cancellation of the two terms linear in \(\mathcal A/D\).
\end{proof}
For the forcing there is also an endpoint expression free of separate exterior sums. With \(b_v=v_t+v v_x\), one has
\begin{equation}
 K-1=\frac{D^2}{g}\{b_v(t,x_+)-b_v(t,x_-)\}.
 \label{eq:acceleration}
\end{equation}
Differentiate the two velocity evaluations along their actual paths. The time derivative of \(v\) is \(\partial_x(J_t/J)\). The unnormalized double-source formula below computes \(J_t\) by inserting \(u^2+u'^2\) in the same integral. A signed bound for this difference could therefore feed the parent's joint forcing barrier directly. A bound on \(\mathcal A\) alone does not bound its time derivative.

\subsection{A positive source for spatial curvature}
Take independent copies \(U,U'\) under \(\E_t\), and use the probability weight \((U+U')/(2a)\) on their joint law. Write \(\E_*\) for this weighted expectation and \(\Delta=U-U'\). Symmetrizing \(\overline B B_x\) proves
\begin{align}
 j(t,x)&=\E_*\cos(x\Delta),&
 w&=\E_*\Delta^2
 =\frac{\E U^3-a\E U^2}{a}=\frac{\dot a}{a},&
 m_4^*&=\E_*\Delta^4 . \label{eq:difference}
\end{align}
The transform \(j\) is even and entire, and \(j(0)=1\). The small frequency in this integral is the difference of the two source frequencies. In particular the quadratic term of \(\log j\) is exactly \(-wx^2/2\).

\begin{proposition}[A finite two-sided gap estimate]\label{prop:finite}
At each time in a common counted chart with endpoints in \([0,X]\), put
\[
 b_* =1-X^2 w/2,\qquad
 l=w-X^2m_4^*/2,\qquad u=w/b_*+X^2w^2/b_*^2 .
\]
If \(b_*>0,\ l\ge0\), then \(l\le-\mathcal A\le u\).
Consequently for \(t\le T\) with these inputs throughout,
\begin{equation}
 D(T)\exp\!\left(2\int_t^T l(r)\dd r\right)
 \ \le D(t)\le\
 D(T)\exp\!\left(2\int_t^T u(r)\dd r\right).
 \label{eq:slab}
\end{equation}
\end{proposition}
\begin{proof}
On \(0\le x\le X\) one has \(b_*\le j\le1\), \(|j'|\le xw\), and
\(-w\le j''\le-w+x^2m_4^*/2\), using \(\cos y\ge1-y^2/2\).
Insert these signed bounds into \(j''/j-(j'/j)^2\), then integrate \eqref{eq:velocity}. The condition \(l\ge0\) permits division of the negative upper bound by \(j\le1\).
\end{proof}

\subsection{The higher exterior moments from one deflated field}
At a fixed counted pair put \(Q(z)=(z-c)^2-D\) and
\(R=\partial_z\log(H/Q)\), analytic near the target segment in a cleared chart. Set \(r^2=D+\xi\), with \(\xi\) small enough that the probes \(c\pm r\) stay in that chart and \(|\xi|<\inf_{\rho\ {\rm ext}}|Q_\rho|\). The parent's complete rational traces \cite[\S3]{AGC} then give
\begin{align}
 \sum_{n\ge0}S_{n+1}\xi^n
 &=-\frac{R(c+r)-R(c-r)}{2r}, \label{eq:Sgen}\\
 \sum_{n\ge0}T_{n+2}\xi^n
 &=-\frac{R(c+r)+R(c-r)-R(c+\sqrt D)-R(c-\sqrt D)}{2\xi},\label{eq:Tgen}
\end{align}
where \(T_n=\sum_{\rho\ {\rm ext}}(\rho-c)Q_\rho^{-n}\) for \(n\ge2\). Quotients take their removable values where necessary. In fact an exterior pole at \(c+d\) contributes \((d^2-r^2)^{-1}\) to \eqref{eq:Sgen}. The symmetric sum in \eqref{eq:Tgen} before subtraction contributes \(-2d/(d^2-r^2)\); its subtraction at \(\xi=0\) removes the first-order ambiguity. The resulting series converge absolutely by the parent's reciprocal-square tail bound. These series generate the indicated diagonal and midpoint-weighted traces at a fixed pair. Section~\ref{sec:exterior} uses both endpoint exponents independently. Controlling its coefficients at lower time requires a positive radius and a bound on the same deflated field.
